\subsubsection{Chování projektilu po vystřelení}

Skript \texttt{BulletStripeShot.cs} se nachází na instanci projektilu a aktivuje se až po \newline vystřelení. V \texttt{OnTriggerEnter2D()} se kontroluje, s jakým objektem projektil kolidoval.

Při vstupu do trigger collideru bowlingové koule (která má tag \texttt{Ball}) se volá metoda \texttt{Stun()} na kouli. To způsobí, že koule se zpomalí a poskytne hráči čas na únik. Na rozdíl od tradičních her, kde střelba eliminuje nepřítele, v naší hře střelba pouze dočasně neutralizuje hrozbu.

Při srážce s prostředím (Collider2D na ground vrstvě) se projektil okamžitě zničí voláním \texttt{Destroy(gameObject)}. Toto je implementováno pomocí layer comparison nebo tag checking, aby se rozlišilo mezi prostředím a jinými objekty.

\begin{lstlisting}[language=csharp,caption={Listing 13: Skript BulletStripeShot}]
public class BulletStripeShot : MonoBehaviour
{
    public bool stun = true;
    public LayerMask groundLayer;

    private void OnTriggerEnter2D(Collider2D collision)
    {
        if (collision.CompareTag("Ball"))
        {
            BallMovement ball = collision.GetComponent<BallMovement>();
            if (ball != null)
                ball.Stun(stun);

            Destroy(gameObject);
        }
    }

    private void OnCollisionEnter2D(Collision2D collision)
    {
        if (((1 << collision.gameObject.layer) & groundLayer) != 0)
            Destroy(gameObject);
    }
}
\end{lstlisting}
